\chapter{Concept Localisation via Contrastive Residual-Stream Analysis}
\label{sec:concept_localisation}

\section{Contrastive Pair Construction}
\label{sec:concept_pairs}

% TODO: Explain contrastive prompts, and contrastive datasets, and why using constrastive pairs would help us find abstract concepts related to mathematical algos.
% TODO: Explain the datasets (carry, gcd, residue class), and write a table, and maybe plot the datasets (and some modular plots on the gcd/residue class datasets, to see
% how data is distributed). Explain templates too (and what they mean, T0- mathy, T1- wordy, T2-question like).
...
The templates differ in whitespace and surface phrasing. Repeating the extraction across all
three provides a check on whether the extracted direction reflects the model's carry
representation or a pattern tied to one specific tokenisation.

\section{Layer-Wise Concept Deltas}
\label{sec:concept_delta}

Explain what deltas means, and how it is mathematically constructed in this project. 
...
like
Let $\mathcal{D}_+$ and $\mathcal{D}_-$ denote the sets of positive and negative prompts
respectively. For each layer $l \in \{0, \ldots, L-1\}$, the concept direction is defined
as the mean difference of residual-stream vectors at the anchor position,
\begin{equation}
\delta_l \;=\; \mathbb{E}_{x \in \mathcal{D}_+}\!\left[ h_{l,t^*}(x) \right]
-\; \mathbb{E}_{x \in \mathcal{D}_-}\!\left[ h_{l,t^*}(x) \right],
\label{eq:carry_delta}
\end{equation}
where $h_{l,t}(x) \in \mathbb{R}^d$ is the residual-stream vector at layer $l$ and position $t$
for prompt $x$. Under the linear representation hypothesis~\citep{park2023linear}, $\delta_l$ approximates the axis along which
concept is encoded at depth $l$, provided the concept is linearly separable and the matched-pair
averaging suppresses nuisance variation. 

% TODO: Explain the anchors, and what they mean. Also, explain the null hypothesis and how it helps us distinguish between a genuine concept direction and a spurious difference, 
% and thus which anchors might provide something for the concept (high separation) vs the algorithm itself (low separation).
% Explain causal scores and overlay. 
% Explain how deltas analysis can be used to identify phase transitions, and how it can be used to identify the effective steering window, and what spiky layers
% mean, and how they can help isolate features that might encode concept-specific (or even algorithmic-specific)information. And how important is that
% as an initial proxy for mechanistic localisation, and how it can be used to guide further analysis and interventions.

\section{Sharpness and Concept Localisation}
\label{sec:sharpness}

% TODO: Explain the norm trajectory, and how it can be used to identify phase transitions, and how the sharpness index is defined, and what it means. 
% Explain what metrics we used to select top-k anchors, adn what that means. 
% Also, explain the inter-layer alignment and what it means.
% Explain the delta projection onto feature W_dec vectors, and what it means.
% also explain the the inter-layer cosine alignment, and layers cosine heatmaps, and what they mean. Also run analysis on how heatmaps at specific anchors relate to
% norm trajectories and phases at those anchors. 

\section{Feature-Space Interpretation}
\label{sec:feature_projection}

% Replace the explanation bellow to correspond to current delta projections onto trasncoder features (using W_dec instead), and what it means.

In the dictionary-learning framework applied to transformer internals, a \textit{feature} is a
learned direction $\phi_f \in \mathbb{R}^d$ hypothesised to correspond to a single interpretable
concept or computation, where $d$ is the residual-stream width of the transformer.
Because transformers represent more independent concepts than they have
dimensions, they are believed to superpose many approximately orthogonal feature directions within
the residual stream, with each activation being a sparse signed combination of active
features~\citep{lindsey2025biology}. A transcoder operationalises this for MLP sub-layers as
follows: an encoder $W_l^{\mathrm{enc}} \in \mathbb{R}^{F_l \times d}$ maps the pre-MLP residual-stream
vector to a sparse bottleneck of feature activations, and a decoder $W_l^{\mathrm{dec}} \in
\mathbb{R}^{d \times F_l}$ reconstructs the MLP output as a sum of active feature effects.
Each encoder row $(W_l^{\mathrm{enc}})_f$ is therefore a \textit{detection direction} in $\mathbb{R}^d$.

The delta vectors $\delta_l$ are dense in $\mathbb{R}^d$ and do not directly identify which
model components encode carry at each layer. The transcoder architecture provides a more interpretable way to
decompose them sparsely. Each transcoder at layer $l$ learns an encoder weight matrix
$W_l^{\mathrm{enc}} \in \mathbb{R}^{F_l \times d}$, where $F_l$ is the number of learned
features and $d$ the residual-stream dimension. The $f$-th row $(W_l^{\mathrm{enc}})_f$ is
the direction to which feature $f$ responds most strongly.

The cosine alignment between the carry delta and each encoder direction,
\begin{equation}
c_{l,f} \;=\;
\frac{(W_l^{\mathrm{enc}})_f \cdot \delta_l}
     {\lVert (W_l^{\mathrm{enc}})_f \rVert \cdot \lVert \delta_l \rVert},
\label{eq:feature_cos}
\end{equation}
identifies the transcoder features most aligned with the concept delta at that layer. Features
with $\lvert c_{l,f} \rvert$ close to one are the monosemantic components whose activation
pattern best accounts for the carry versus no-carry difference. This projection connects the
dense, direction-based view of representation engineering to the sparse, feature-based view of
dictionary learning. The same concept direction can thus be understood both as a point in
residual-stream space and as a weighted combination of interpretable transcoder features,
bridging two perspectives that the broader interpretability programme has pursued in parallel.

\section{Template Robustness}
\label{sec:template_consistency}

% Explain template robustness, and how it can be used to check whether the extracted direction is a genuine property of the model's internal computation rather than an artefact of any particular surface form.
% Also, run an analysis on per-template metrics and results, to understand how template structure drives different outcomes in terms of 
% deltas (are phases the same but deltas different? yes/no and why), magnitudes, causal scores, heatmaps (more diffuse or more localized). correlate template types to prior analysis like norm trajectories and phase transitions and heatmaps.

...

Per-template deltas $\delta_l^{(T)}$ are computed separately for each template
$T \in \{T0, T1, T2\}$, and their pairwise cosine similarities are measured at each layer.
This provides a direct check on whether the extracted direction is a stable property of the
model's carry representation. Consistently high cross-template alignment at the layers where
$\lVert \delta_l \rVert$ peaks would confirm that the same direction emerges regardless of
surface form. Divergence at early layers would indicate that template-specific variation
persists until mid-network processing assimilates the input into a common representational
format, providing additional evidence about the depth at which template-invariant concept
encoding begins. 

\section{Causal Validation and Null Hypothesis Testing}
\label{sec:causal_validation}

%  Explain these in detail here
The delta vectors and sharpness metrics of the preceding sections characterise the
\textit{representational} structure of each concept, quantifying where a separable direction
exists in residual-stream space. They do not, however, establish that the extracted
representation is causally upstream of the model's output. A direction may be strongly
present at a given layer without being mechanistically connected to the prediction. Two
complementary interventions test this causal claim directly, and are applied to every concept
in the suite.

\subsection{Activation Patching}
\label{sec:act_patching}

Activation patching~\citep{meng2022locating} isolates the causal contribution of a
specific layer by replacing a single intermediate representation and measuring the
resulting change in output. For each contrastive pair $(x^+, x^-)$ and each layer $l$,
the positive prompt's residual-stream vector at the anchor position $t^*$ is transplanted
into the negative prompt's forward pass. The score is the induced change in the
logit margin at the final prediction position,
\begin{equation}
S_l^{\mathrm{patch}}
\;=\;
\Bigl[
  f_{\mathrm{pos}}\!\left(\tilde{x}^-_{[l]}\right) - f_{\mathrm{neg}}\!\left(\tilde{x}^-_{[l]}\right)
\Bigr]
-
\Bigl[
  f_{\mathrm{pos}}\!\left(x^-\right) - f_{\mathrm{neg}}\!\left(x^-\right)
\Bigr],
\label{eq:act_patching}
\end{equation}
where $\tilde{x}^-_{[l]}$ denotes the negative prompt with $h_{l,t^*}(x^-)$ replaced
by $h_{l,t^*}(x^+)$, and $f_{\mathrm{pos}}$, $f_{\mathrm{neg}}$ are the logits for
the positive and negative answer tokens at the last position. A large positive $S_l^{\mathrm{patch}}$
indicates that the layer-$l$ concept representation, when inserted into the negative
context, is sufficient to shift the model's output toward the positive answer. The metric
therefore directly quantifies the \textit{causal sufficiency} of the layer-$l$ encoding
for driving the correct response.

\subsection{Gradient Dot Delta}
\label{sec:grad_dot_delta}

Activation patching requires a full forward pass for each patched layer, making it
the more expensive of the two methods. A cheaper first-order approximation is obtained
by differentiating the output margin with respect to the residual stream at the negative
prompt's operating point. For each pair and each layer $l$, a single forward pass on
$x^-$ is followed by a single backward pass to compute the gradient of the margin
with respect to $h_{l,t^*}(x^-)$. The score is the inner product of this gradient with
the precomputed concept delta,
\begin{equation}
S_l^{\mathrm{grad}}
\;=\;
\nabla_{h_{l,t^*}} m\!\left(x^-\right) \cdot \delta_l,
\label{eq:grad_dot_delta}
\end{equation}
where $m(x) = f_{\mathrm{pos}}(x) - f_{\mathrm{neg}}(x)$ is the logit margin evaluated
at prompt $x$, and the gradient is computed at the negative prompt's residual stream.

The relationship between $S_l^{\mathrm{grad}}$ and $S_l^{\mathrm{patch}}$ is made
precise by a first-order Taylor expansion of the patched margin around the unpatched
activation,
\begin{equation}
m\!\left(\tilde{x}^-_{[l]}\right)
\;\approx\;
m\!\left(x^-\right)
+ \nabla_{h_{l,t^*}} m\!\left(x^-\right)
  \cdot \bigl(h_{l,t^*}(x^+) - h_{l,t^*}(x^-)\bigr).
\label{eq:taylor_patch}
\end{equation}
Taking the expectation over pairs and using $\delta_l = \mathbb{E}[h_{l,t^*}(x^+)] - \mathbb{E}[h_{l,t^*}(x^-)]$ gives
$\mathbb{E}[S_l^{\mathrm{patch}}] \approx S_l^{\mathrm{grad}}$ to first order, provided the gradient
varies little across the pair distribution. The gradient-dot-delta score is therefore the
linearised, single-pass analogue of activation patching; agreement between the two methods
at any given layer constitutes evidence that the output function is approximately
linear in the relevant neighbourhood of residual-stream space. Systematic divergence
indicates nonlinearity or higher-order interactions that the gradient cannot capture.
Both methods are anchored at the same position $t^*$ used for delta extraction throughout
this chapter.

%  Add Null hypothesis explanation

\section{Results}


% TODO LATER...